\section{Conclusion} \label{section: conclusion}

This work started from a concrete question: whether learned KV-cache eviction, which the recent literature trains with bespoke implementations such as Apple's (KVP, \cref{subsec: kvp}), could be reformulated so as to be expressible as a standard Gymnasium environment and trainable with Stable-Baselines3's masked PPO. The central idea was to turn eviction, which in KVP is a single-pass ranking of the whole cache, into a \emph{sequence} of discrete decisions, and moreover a dynamic one, as the model generates tokens. That reformulation did turn out to be trainable, which is already a step beyond the state of the art.

The central empirical result is that the answer to that question depends on the dataset, not on the method. On GSM8K, no axis we attacked (representation, architecture, initialisation, credit assignment, reward density) took the policy beyond parity with \texttt{kv\_norm}. But measured with a future-attention oracle, which no online policy can reach, the margin separating \texttt{kv\_norm} from optimal eviction on GSM8K is barely $+0.07$; on retrieval tasks, real ($+0.09$ to $+0.19$ depending on the compression) or synthetic ($+0.43$), that margin is several times larger. GSM8K, precisely the dataset chosen to test eviction \emph{during} generation and not only over a fixed context, turned out to be the worst possible case for this problem. Once a dataset with signal was located, our own formulation, online, with the causal dense reward and a configuration that stabilises training, shows there a clear shift above \texttt{kv\_norm} (though not yet stable, and resting for now on a single seed).

This also lets us revisit the design decision discussed in \cref{subsec: ngc}: keeping the language model frozen and training an external, lightweight policy, instead of training the whole model as NGC does. A policy that only sees $K$ and $V$ has no access to everything the model ``knows'' internally about a token's future utility. An approach like NGC, which retrains the model while deciding what to forget, could exploit that internal information. The contrast between the two approaches is clearer after this work than before: it is not only a question of computational cost, but of what information is available depending on where the policy sits. With more time and compute the two worlds could be joined into a method unifying the most advanced techniques in the state of the art.

\subsection{Contributions}

This work makes two main contributions. First, a sequential reformulation of learned eviction that makes it expressible as a Gymnasium environment and trainable with standard masked PPO, together with the online eviction environment over a frozen model that implements it. And second, the central finding: the conditions under which learned eviction can beat a heuristic, together with preliminary evidence that our own online reformulation can exploit them.

\subsection{Limitations and future work}

The work has three limitations that bound its conclusions. The first is one of scope: every experiment was run on a single model (Qwen2.5-1.5B-Instruct) and moderate cache budgets, and absolute accuracy can depend on the attention backend used in each run, so only the paired contrasts within a single run are comparable with one another. The second concerns the causal dense reward: it rewards the policy for resembling, step by step, a model that never evicts, but that is not the same as directly rewarding the correctness of the final answer; an eviction can shift the token distribution only slightly and still change the answer, or shift it a great deal without affecting it. That decoupling between what the reward measures and what ultimately matters is, in part, what keeps making training difficult even when there is signal to exploit. The third is that the positive result of our online method is still preliminary: it rests on a single hyperparameter configuration and, for now, on a single seed, and it needs confirming before being treated as a firm conclusion.

Several directions follow from this. The most immediate is to confirm the online result with more seeds and longer training, and to incorporate PPO stabilisation techniques to turn the observed shift into sustained convergence. In parallel, it would be valuable to take the environment to real long-context datasets, where the learnable margin measured by the oracle is larger, and to explore per-layer credit assignment based on correctness and not only on distributional divergence. Finally, the contrast with approaches that train the model and the policy jointly, such as NGC, remains a natural comparison for understanding how much of the difficulty of this work comes specifically from learning to forget over a model that cannot adapt to that constraint.
